% Transcription — fonds Alexandre Grothendieck, shelfmark 71, batch 3 (pages 41-47,
% the end of the folder). Established from the facsimile
% archives/batches/71/batch-03.pdf (a mirror of
% https://grothendieck.umontpellier.fr/71.pdf), per the skill
% transcribe-grothendieck. The batch file carries no cover sheet: PDF page k is
% archivists' page 40+k; the pencilled numbers bottom-left (« 41 » ... « 47 »)
% agree.
%
% Pass: Opus 5.5 (claude-opus-5-5), 2026-09-24 — first pass, unchecked against the pages by a human.
% Revised 2026-09-24 (Opus 5.5), by the user's decision of that day: the
% mémoire is summarised, not recomposed (see below).
%
% What the pages are. All seven are the last leaves of « un mémoire d'une
% autre main, signé et daté du 15 juillet 1976 » — a fair copy in blue ink,
% paginated by its author in boxed Roman numerals XIV (page 41) to XIX (page
% 47). None is a working leaf of Grothendieck's. Following folder 100's rule
% for other people's text (header of transcripts/100/batch-02.fr.tex), by the
% user's decision of 2026-09-24, each mémoire page keeps its \page and is
% summarised in one French \note — its page number in the mémoire and the
% lemmas it states, one sentence each, no proofs — and is not recomposed.
% The marks of another hand are given in full as \marginal, each tied to the
% words of the mémoire it answers.
%
% Summarised: pages 41, 42, 43, 44, 46, 47.
% Dropped: page 45, the verso of the leaf of page 44 — two pencil drawings of
% non-convex polygons (vertices s', s, s'') and, upside down at its foot, a
% false start « ment [s,s'] ∩ » in the mémoire's ink. No \page.
%
% The annotations. On pages 46-47 a second hand, in pencil, has annotated the
% copy: an insertion into Lemme XX, question marks, « ! », « so what »,
% « ou bien, », a circled s'', an underlining, and on page 47 a slanted note
% read only in part (« faux, si on ... choisit ... »). Whose hand this is is
% not established: it is not the mémoire's, and it may be a reader's —
% possibly his — but the file does not assert it. On page 43 a pencil circle
% surrounds the mémoire's own insertion « XIV ». Pages 41, 42 and 44 carry no
% annotation in words (page 42 has a few wordless vertical strokes in ink in
% both margins).
%
% What the mémoire does here. The end of a proof of the Jordan theorem for
% polygons: Lemmes XVI-XVII (two components U, V of R^2 - Γ; Théorème 1),
% § 4 « Preuve du théorème 2 » by induction on card S(Γ) - 2 with the
% polygons Λ^1, Λ^2 cut along [s,s'], Lemmes XVIII-XXI (relative position of
% Λ^1, Λ^2; good vertices), and the conclusion of Théorème 2.
%
% The dating: the folder's own, which the mémoire's date on page 47 bears
% out. The group « Géométrie et topologie combinatoire » (69-89) holds it.
\documentclass[11pt,a4paper]{article}
% Le préambule commun à toutes les transcriptions du fonds Grothendieck.
%
% Il tient en une idée : **l'incertitude se compose, elle ne se résout pas.**
% Une transcription qui lisse un mot douteux a détruit la seule chose pour
% laquelle on la faisait — un lecteur doit pouvoir savoir, à chaque endroit,
% ce qui était sur le feuillet et ce qui a été ajouté. Les six macros qui
% suivent sont donc l'appareil critique, et non de la décoration.
%
% Les mêmes commandes sont reconnues par `scripts/render.mjs`, qui produit la
% vue de lecture du site. Une macro ajoutée ici doit l'être aussi là-bas,
% faute de quoi le rendu échouera bruyamment — ce qui est voulu : un rendu
% silencieusement faux est pire qu'un rendu absent.

\ProvidesPackage{grothendieck}[2026/08/08 Transcriptions du fonds Grothendieck]

\usepackage{amsmath,amssymb,amsthm}
\usepackage{mathrsfs}
\usepackage{stmaryrd}
% Les diagrammes commutatifs sont partout dans ce fonds : c'est la langue
% même de la géométrie algébrique de Grothendieck, et une transcription qui
% les rendrait en prose ne transcrirait rien.
\usepackage{tikz-cd}
% Français par défaut — c'est la langue des feuillets — mais l'anglais est
% chargé aussi : la traduction et le résumé sont des documents anglais, et la
% typographie française (espaces avant les deux-points) y serait une faute.
% Ces documents basculent par \selectlanguage{english} en tête de corps.
\usepackage[english,french]{babel}
\usepackage{fontspec}
\usepackage{geometry}
\geometry{a4paper,margin=25mm,marginparwidth=28mm}
\usepackage{marginnote}
\usepackage{xcolor}
% ulem, not soul. `soul`'s \st and \ul are built on a character-by-character
% scan that XeLaTeX's font machinery breaks: the document fails with an
% undefined control sequence rather than degrading. ulem does the same two jobs
% and is Unicode-engine safe. `normalem` stops it hijacking \emph, which the
% transcriptions use for Grothendieck's own emphasis.
\usepackage[normalem]{ulem}
\usepackage{hyperref}
\hypersetup{colorlinks=true,linkcolor=black,urlcolor=black,pdfborder={0 0 0}}

% --- Métadonnées du lot -----------------------------------------------------
% Reprises telles quelles par le script de rendu, qui les affiche en tête de
% la vue de lecture. Elles fixent ce que le fichier prétend couvrir : sans
% elles, une transcription flotte sans point d'attache dans un fonds de seize
% mille feuillets.

\newcommand{\folder}[1]{\gdef\grothFolder{#1}}
\newcommand{\batch}[1]{\gdef\grothBatch{#1}}
\newcommand{\leaves}[2]{\gdef\grothFirst{#1}\gdef\grothLast{#2}}
\newcommand{\foldertitle}[1]{\gdef\grothFoldertitle{#1}}
\gdef\grothFolder{?}\gdef\grothBatch{?}\gdef\grothFirst{?}\gdef\grothLast{?}\gdef\grothFoldertitle{}

% --- Le résumé --------------------------------------------------------------
%
% Il ouvre la lecture modernisée, sous le titre « Résumé », et il est composé
% en petit. Ce n'est pas un ornement : c'est le seul passage du document qui
% ne soit pas des mathématiques, et un lecteur doit pouvoir le reconnaître
% avant de l'avoir lu — puis le sauter s'il connaît déjà le sujet, ou s'y
% arrêter s'il ne le connaît pas. Le filet qui le suit marque la reprise du
% corps.

\newenvironment{resume}
  {\small\setlength{\parskip}{0.45em}}
  {\par\medskip\hrule\medskip}

% --- Mots-clés ---------------------------------------------------------------
%
% En anglais, en clôture du résumé de la lecture modernisée : le vocabulaire
% moderne sous lequel chercher ce que les feuillets construisent. Le manifeste
% du site (`npm run manifest`) les extrait pour étiqueter la cote entière —
% cette ligne est la source unique des tags, il n'y a pas de fichier à côté.
\newcommand{\keywords}[1]{\par\smallskip\noindent
  {\footnotesize\textsc{Keywords} --- \textit{#1}}\par}

% --- L'appareil critique ----------------------------------------------------

% Le numéro de feuillet, en marge. C'est l'ancre qui permet de régler un
% désaccord sur une formule en regardant *un* feuillet plutôt qu'une chemise
% de six cents. Le site s'en sert aussi pour tourner les pages du fac-similé
% quand on fait défiler la transcription.
\newcommand{\leaf}[1]{%
  \marginnote{\footnotesize\color{gray!70}\textsf{#1}}%
  \label{leaf:#1}\ignorespaces}

% Illisible : on ne devine pas. Un mot inventé qui se lit comme les autres est
% le pire résultat possible.
\newcommand{\ill}{\textcolor{red!60!black}{[\,\ldots\,]}}

% Lecture douteuse : le mot est donné, et signalé comme incertain.
\newcommand{\uncertain}[1]{\uline{#1}}

% Ajout de l'éditeur — une lettre manquante, un mot sous-entendu.
\newcommand{\add}[1]{\textcolor{blue!50!black}{[#1]}}

% Biffé par Grothendieck lui-même. Ce qu'il a rayé fait partie du document :
% une rature dit souvent où la pensée a changé de direction.
\newcommand{\struck}[1]{\sout{#1}}

% Note du transcripteur, sur l'état matériel du feuillet.
\newcommand{\note}[1]{\par\smallskip{\small\color{gray!60!black}\textit{[#1]}}\par\smallskip}

% Note marginale de Grothendieck, distincte du corps du texte.
\newcommand{\marginal}[1]{\par\smallskip\noindent\hspace{1em}%
  {\small\color{gray!60!black}#1}\par\smallskip}

% --- Titre ------------------------------------------------------------------

\AtBeginDocument{%
  \begin{center}
    {\large\bfseries Fonds Alexandre Grothendieck --- cote n\textsuperscript{o}~\grothFolder}\\[2pt]
    {\itshape\grothFoldertitle}\\[4pt]
    {\small feuillets \grothFirst--\grothLast\ (lot \grothBatch)}\\[2pt]
    {\footnotesize\color{gray!60!black}Transcription établie d'après le fac-similé numérisé
     par l'Université de Montpellier.\\
     Les lectures incertaines sont soulignées, les illisibles notées [\,\ldots\,],
     les ajouts entre crochets.}
  \end{center}
  \vspace{1em}}

\endinput


\folder{71}
\batch{3}
\pages{41}{47}
\dating{[à partir de 1976]}
\watermark{Édition de démonstration}
\foldertitle{Théorème de Jordan : notes manuscrites (1976, s.d.).}

\begin{document}

\page{41}
\note{Page XIV d'un mémoire d'une autre main, signé et daté du 15 juillet
1976, résumée ici et non recomposée. \emph{Lemme} XVI : a) $\mathbb{R}^2 -
\Gamma = U \cup V$ ; b) $U$ et $V$ sont des ouverts connexes, $U$
relativement compact, $V$ non relativement compact. La page démontre a) et
commence la connexité de $U = \operatorname{int}\Lambda^1 - R_{\Lambda^2}$
(où $R_{\Lambda^2} = \operatorname{int}\Lambda^2 \cup \Lambda^2$), par des
arcs de polygone $\Lambda_{xy}$ et la Remarque 1. Aucune annotation.}

\page{42}
\note{Page XV du mémoire. Fin de la démonstration du lemme XVI : connexité
de $U$ (voisinage $V$ de $\Gamma_{x'y'}$, lemme I), puis de $V =
\operatorname{ext}\Lambda^1 \cup \operatorname{int}\Lambda^2 \cup\,
]s,s'[$ ; $U$ et $V$ sont ouverts comme composantes connexes de
$\mathbb{R}^2 - \Gamma$, $\Gamma$ étant compact. Quelques traits verticaux à
l'encre, sans mot, dans les deux marges.}

\page{43}
\note{Page XVI du mémoire. \emph{Lemme} XVII : si $x \in \Gamma$ et $D_x$
est un disque de centre $x$ qui ne coupe que les arêtes auxquelles $x$
appartient, alors $D_x \cap U \neq \emptyset$ et $D_x \cap V \neq
\emptyset$. Les lemmes XVI et XVII donnent le théorème 1. § 4, « Preuve du
théorème 2 » (tout polygone est un $J$-polygone), par récurrence sur
$\operatorname{card} S(\Gamma) - 2$ : $s'$ visible de $s$ de l'intérieur
(lemme XI), les arcs $\Gamma^1_{ss'}$, $\Gamma^2_{ss'}$, et les polygones
$\Lambda^1$, $\Lambda^2$ obtenus en leur recollant $[s,s']$, qui ont moins
de sommets que $\Gamma$. Le renvoi « par le lemme XIV », ajouté par
l'auteur, est entouré au crayon.}

\page{44}
\note{Page XVII du mémoire. \emph{Lemme} XVIII : $\Lambda^2 - \{s,s'\}
\subset \operatorname{ext}\Lambda^1$ et $\Lambda^1 - \{s,s'\} \subset
\operatorname{ext}\Lambda^2$. \emph{Lemme} XIX :
$\operatorname{int}\Lambda^1 \subset \operatorname{ext}\Lambda^2$,
$\operatorname{ext}\Lambda^1 \supset \operatorname{int}\Lambda^2$. La page
donne la démonstration du lemme XVIII et le début de celle du lemme XIX.
Aucune annotation.}

\page{46}
\note{Page XVIII du mémoire. Fin de la démonstration du lemme XIX.
\emph{Lemme} XX : il y a deux bons sommets de $\Gamma$ tels que les
intérieurs des triangles correspondants soient contenus dans
$\operatorname{int}\Gamma$. La démonstration applique l'hypothèse de
récurrence à $\Lambda^1$ et $\Lambda^2$, puis, si $t_1 = s$, distingue
$]s'',s'[\, \cap\, \Gamma = \emptyset$ et $]s'',s'[\, \cap\, S(\Gamma)
\neq \emptyset$ en construisant un polygone $\Lambda^{1\prime}$ ; elle se
poursuit page 47. Les annotations au crayon qui suivent sont d'une autre
main que le mémoire ; la leur n'est pas établie.}

À « Il y a deux bons sommets de $\Gamma$, \struck{au moins} tels que »,
inséré par un renvoi après « de $\Gamma$ » (parenthèse non fermée) :
\marginal{(non consécutifs si $n$ = \uncertain{nb de côtés} de $\Gamma
\geq 4$,}

À « il y a un bon sommet $t_1$ dans $\Lambda^1$ », où « bon sommet » est
souligné au crayon, deux ou trois lettres en biais dans la marge :
\marginal{\ill{}}

À « Deux cas peuvent se produire $]s'',s'[\, \cap\, \Gamma = \emptyset$ »,
le premier $s''$ étant cerclé (il n'est pas défini dans le texte) :
\marginal{?}

À « et dans ce cas $]s'',s'[\, \cap\, S(\Gamma) \neq \emptyset$ », un signe
griffonné : \marginal{\ill{}}

À « obtenu en recollant selon $\{s', s''\}$ l'arc de polygone $[s',s'']$ » :
\marginal{!}

À « l'ensemble des arêtes de $\Lambda^{1\prime}$ contenues dans $A(\Gamma)$
est de cardinal inférieur » : \marginal{?}

À « au cardinal de l'ensemble des arêtes de $\Lambda^1$ contenues dans
$A(\Gamma)$ » : \marginal{so what}

\page{47}
\note{Page XIX, dernière page du mémoire. Fin de la démonstration du lemme
XX, par une descente sur le nombre d'arêtes de $\Lambda^{1(n)}$ contenues
dans $A(\Gamma)$, en deux cas a) et b). \emph{Lemme} XXI : si $\Gamma$ est un
polygone non convexe, il existe un bon sommet de $\Gamma$ tel que
l'intérieur du triangle correspondant est contenu dans
$\operatorname{ext}\Gamma$. Les lemmes XX et XXI, joints au fait qu'un
polygone convexe a au moins trois bons sommets, donnent le théorème 2. Le
mémoire se clôt sur « le 15 juillet 1976, Montpellier » et une signature.
Annotations au crayon de la même main qu'à la page 46.}

Devant « b) il existe $n$ tel que $\operatorname{card}(A(\Lambda^{1(n)})
\cap A(\Gamma)) = 2$ » : \marginal{ou bien,}

À « il est clair que $\operatorname{int}\Lambda^2 \subset
\operatorname{ext}\Gamma$ car $V = \operatorname{ext}\Lambda^1 \cup
\operatorname{int}\Lambda^2 \cup\, ]s,s'[$, et $V$ est non-relativement
compact », quatre lignes en biais le long de la marge, lues en partie :
\marginal{\uncertain{faux, si on} \ill{} \uncertain{choisit} \ill{} \ill{}
\ill{}}

À « les deux arêtes de $\Lambda^2$ auxquelles $s$ appartient », un mot en
biais, souligné d'un long trait : \marginal{\ill{}}

\end{document}
